% Modernised reading — fonds Alexandre Grothendieck, shelfmark 156-1, pages 1-26.
%
% This is an INTERPRETATION, not a transcription. Everything the critical
% apparatus carried moves into footnotes. In any dispute of substance the
% transcription governs: whoever cites Grothendieck must cite batch-01.fr.tex
% (pages 1-20) and batch-02.fr.tex (pages 21-26).
%
% Pass: Opus 5.5 (claude-opus-5-5), 2026-10-10 — first pass, unchecked against
% the pages by a human. The folder was read whole: two batches, pages 1-26, all
% mathematics, dated 5 June 1986 (page 1), 7 June (margin of page 3, page 23).
% Not measured against the Opus 5 reference reading of folder 115.
%
% What the folder is. Chapter I (« GF I ») of Vers une géométrie des formes: a
% purely combinatorial axiomatics of one-dimensional « formes » — a set of
% « lieux » and a set of « segments », each with its set of lieux and its two
% extremities — in three successive moutures (pages 1-2, 3-4, 5-26), kept apart
% on the spine. The third mouture reaches axioms F1-F6, the « modèle-segment »
% (= a bounded, dense, locally filtering order, page 13), orientation of
% regular connected models (pages 15-17), and the « réseaux » (page 23), whose
% reduced case is a simple graph. Chapter II (folder 156-2), begun on 6 June,
% realises the same notions topologically; its page 9 lists for the lieux of a
% faille exactly the order properties of page 13 here. Page 23 here is dated
% 7 June, so chapter I was still being written after chapter II had begun.
%
% Notation fixed for the whole document. L the set of lieux, S the set of
% segments, I~ the set of lieux on I, ∂I ∈ P_2(L) its extremities, I° = I~ ∖ ∂I.
% From Cor. 4 (page 6) on, a segment is identified with I~ and written I, as
% the page does from page 21. I_{x,z} the sub-segment of I with extremities
% x, z (F1). L_r, L_s regular and singular lieux. S_x, R_x, B_x = S_x/R_x the
% segments issued from x, the relation « même branche », the branches;
% br_x(I) the branch of I at x. [x,y] = {z : x ≤ z ≤ y}. ω the set of
% orientations. « Disjoints » for two lieux keeps the page's sense
% (comparable, « emboîtés »), never mere distinctness.
%
% Substantive departures from the pages, each footnoted where it occurs:
% — page 6: R_x is transitive only with a filtering condition at x that F1-F3
%   do not give; the margin of page 6 begins to write it (« Il manque … F'3 »);
%   we state it as F'3, the reconstruction ours.
% — page 7: « on a … l'égalité ! » in F4 is not taken as I = I' ∪ I'': it fails
%   in every model whose order is not total (page 4's margin warns against it);
%   only the inclusion is asserted.
% — page 7, Cor. 1: the sketch needs K with K ∩ I = {x}; F3 gives one meeting
%   only I_{x,u} in {x}. The gap is said, not filled.
% — page 9, Remarque 2: F4 passes to (L', S') only if the glued segment stays in
%   L'; said. F1-F3, F5 pass.
% — page 18: « si x ≤ y, z, ∃ u, x ≤ u ≤ y, z » is trivial as written; read
%   strictly.
% — page 21: « x' < a < a'' » read as ∀ a, ∃ a' < a < a''.
% — page 22: F6 is stated with interiors, consequence a) with closed segments;
%   the step is not on the page; the argument of pages 20 and 25 needs the
%   closed form; said.
% — page 23: a « réseau réduit » (I° = ∅) contradicts F2 as stated on page 5
%   (I° ≠ ∅); réseaux are read with that clause of F2 dropped and the branch
%   axiom in its margin form (u ∈ I ∖ {x}), which is what makes « les axiomes
%   sont vides ! » (page 24) true.
%
% Modern names given and footnoted as ours: ordre dense, ensemble ordonné
% filtrant (directed set), ordre borné, intervalle, relation d'intermédiarité
% (betweenness, Pasch 1882, Huntington, Hilbert), graphe simple, degré d'un
% sommet, torseur des orientations / revêtement d'orientation.
%
% What the folder announces and never establishes: the transitivity of R_x
% from F1-F3; the equivalence F5 ⇔ F'5 (page 8); the theorem of page 15 that
% the orientation relation has exactly two classes (« si c'était faux, il
% faudrait le mettre en axiome »); the proof that the order of the linear case
% is filtering (pages 19-20), which leads to F6 instead; the study of two
% segments with the same extremities (pages 25-26), which stops mid-sentence;
% the equivalence relation on singular lieux announced on page 9 (and posited
% as E1 on page 2); the recursion on dimension n announced on page 1.
\documentclass[11pt,a4paper]{article}
\input{../preamble/grothendieck}

\folder{156-1}
\pages{1}{26}
\dating{1986}
\watermark{Édition de démonstration\\interprétation personnelle de l'œuvre}
\foldertitle{[Chapitre] I. Vers une géométrie des formes (topologiques) : notes manuscrites (05/06/1986) — lecture modernisée du dossier entier}

\begin{document}

\section*{Résumé}

\begin{resume}
Qu'est-ce qu'une courbe, si l'on oublie les nombres réels, les distances et
même les ouverts ? Ces vingt-six feuillets, écrits du 5 au 7 juin 1986,
essaient de le dire avec deux mots seulement : des \emph{lieux}, et des
\emph{segments}. Un segment a deux extrémités et porte des lieux ; un lieu
intérieur à un segment le coupe en deux segments plus petits ; deux segments
qui se touchent bout à bout en un lieu ordinaire se recollent en un seul. Tout
le reste --- points de branchement, branches issues d'un lieu, orientation,
réseaux --- doit sortir de quelques axiomes portant sur ces deux mots.

C'est le premier chapitre d'un texte que Grothendieck intitule « Vers une
géométrie des formes (topologiques) », et la première page annonce un
programme plus vaste : construire par récurrence sur $n$ une géométrie des
formes de dimension $\leq n$. Seule la dimension $1$ est abordée. On peut
penser à la manière dont Pasch et Hilbert décrivent la droite par la
relation « $z$ est entre $x$ et $y$ » ; mais ici le segment n'est pas
supposé totalement ordonné. Le chapitre suivant (dossier 156-2) montrera
pourquoi : les « lieux » d'une bande topologique, ce sont les courbes qui la
coupent en deux, et deux telles courbes peuvent se croiser sans que l'une soit
au-dessus de l'autre.

Les feuillets reprennent trois fois leur départ. La première mouture (pages 1
et 2) se donne trop de choses --- lieux critiques, relations « disjoints »,
« adjacents », « spécialisation » --- et constate aussitôt que certaines
découlent des autres. La deuxième (pages 3 et 4) réduit les données à deux
ensembles et deux applications, et est barrée d'une grande croix. La troisième
(à partir de la page 5) tient : axiomes de divisibilité, des branches, de
recollement, des lieux réguliers. Elle démontre qu'un segment est la même
chose qu'un ensemble ordonné muni d'un plus petit et d'un plus grand élément,
dense et localement filtrant, puis cherche à orienter les modèles connexes
sans branchement. Là, une démonstration bute : un contre-exemple (page 21)
montre que les axiomes ne suffisent pas, et un sixième axiome est ajouté. Le
7 juin, l'objet ainsi défini reçoit un nom, \emph{réseau}, et dans le cas
extrême où les segments n'ont pas de lieux intérieurs, ce n'est rien d'autre
qu'un graphe.

On voit sur ces pages comment un axiome naît chez Grothendieck. Il écrit « Je
présume », « Je ne sais s'il faut supposer », puis tranche deux jours plus
tard dans la marge ; et devant un énoncé qu'il ne sait pas encore démontrer :
« Si c'était faux, il faudrait le mettre en axiome. S'il est vrai, il doit
être immédiat. » L'axiomatique n'est pas posée d'avance : elle est ce que les
exemples réclament, et elle change quand un exemple la défait.

\keywords{betweenness relation, bounded poset, dense order, directed set, locally directed order, axiomatic one-dimensional geometry, branch of a curve, valence, orientation double cover, simple graph, combinatorial graph}
\end{resume}

\section*{Le fil du dossier, et les conventions}

\begin{enumerate}
\item[1.] \emph{Le projet, et la première mouture} (pages 1 et 2).
\item[2.] \emph{Deuxième mouture} (pages 3 et 4).
\item[3.] \emph{Troisième mouture : données, divisibilité} (pages 5 et 6) ---
axiomes F1, F2 et corollaires 1 à 4.
\item[4.] \emph{Branches} (pages 6 et 7) --- axiome F3.
\item[5.] \emph{Lieux réguliers et recollement} (pages 7 et 8) --- F4, F5.
\item[6.] \emph{Restriction à une partie} (pages 8 et 9).
\item[7.] \emph{Le modèle-segment est un biordre} (pages 10 à 14).
\item[8.] \emph{Modèles réguliers connexes, orientation} (pages 14 à 17).
\item[9.] \emph{Le cas linéaire, et la filtration relative} (pages 17 à 21).
\item[10.] \emph{L'axiome F6} (page 22).
\item[11.] \emph{Réseaux} (pages 23 et 24), daté du 7 juin.
\item[12.] \emph{Deux segments consécutifs} (pages 24 à 26), inachevé.
\end{enumerate}

La page 1 porte « GF I » en angle, le titre, « (Juin 86) » et, en tête du
texte, « (5 juin) ». La pagination de l'auteur recommence avec la deuxième
mouture : notre page 3 est sa page 2, et nos pages 4 à 26 sont ses pages 3 à
24. Quand le texte renvoie à « cor.~2 p.~7 », c'est donc notre page 8. Une note
marginale de la page 3 est datée « 7.6. », comme la page 23 : ce chapitre a
été complété le lendemain du jour où commençait le chapitre II (dossier
156-2, daté du 6 juin).

\emph{Trois moutures.} Les pages 1--2, 3--4 et 5--26 sont trois départs
successifs de la même définition. On ne les fusionne pas : chacune est
présentée pour elle-même, et l'on dit, au passage, ce que la suivante reprend
ou abandonne. La numérotation des axiomes n'est pas la même d'une mouture à
l'autre (le F2 de la page 4 devient le F4 de la page 7) ; dans tout le
document, « F$n$ » sans autre précision désigne l'axiome de la troisième
mouture.

\emph{Conventions.} $\mathcal{L}$ est l'ensemble des lieux, $S$ celui des
segments. Pour $I \in S$, $\widetilde{I} \subset \mathcal{L}$ est l'ensemble
des lieux « sur » $I$, $\partial I \in \mathfrak{P}_{2}(\mathcal{L})$ ses deux
extrémités (une partie à deux éléments), $\partial I \subset \widetilde{I}$,
et $I^{\circ} = \widetilde{I} \smallsetminus \partial I$ ses lieux intérieurs.
À partir du corollaire 4 (page 6), qui montre qu'un segment est déterminé par
ses lieux, on identifie $I$ à $\widetilde{I}$ et l'on écrit simplement $I$,
comme le font les pages 21 à 26 ; c'est le seul changement de notation du
document. $I_{x,z}$ est le sous-segment de $I$ d'extrémités $x$ et $z$
(axiome F1). Le mot « disjoints », appliqué à deux lieux, a le sens de la page
6 et de la page 8 du dossier 156-2 : comparables, emboîtés --- jamais
simplement distincts.

\emph{Ce que le dossier annonce et n'établit pas} : la récurrence sur la
dimension promise à la page 1 ; la relation d'équivalence sur les lieux
singuliers annoncée à la page 9 ; la transitivité de la relation « même
branche » à partir des seuls F1 à F3 ; l'équivalence F5 $\Leftrightarrow$
F$'$5 ; le théorème d'orientation des pages 15 à 17 ; le caractère filtrant
de l'ordre dans le cas linéaire, que les pages 19 et 20 cherchent et qui
conduit à F6 ; l'étude de deux segments de mêmes extrémités, interrompue au
milieu d'une phrase (page 26).

\pagerange{1}{2}
\section*{Le projet, et la première mouture (pages 1 et 2)}

Le texte s'ouvre sur une intention : une « approche vers une construction
récurrente (sur l'entier naturel $n$) d'une géométrie des formes de dimension
$\leq n$ ». Une forme de dimension $0$ serait un ensemble discret, dont les
éléments sont les \emph{lieux}\footnote{Une note oblique dans la marge, en
grande partie illisible, ébauche la notion de modèle d'une forme de
dimension~$0$ : un ensemble $\mathcal{L}$ muni d'une relation, lue avec doute
« équivalence », et d'une relation « disjoints ». On n'en tire rien.}. La
récurrence n'est jamais écrite : tout le dossier traite la dimension~$1$.

Un \emph{modèle de dimension $1$}, dans cette première mouture, comporte :
\begin{enumerate}
\item[(1)] un ensemble $\mathcal{L}$ de lieux ;
\item[(2)] une partie $\mathcal{L}_{0}$ de lieux \emph{critiques}, ou \emph{de
branchement} ;
\item[(3)] une relation symétrique « $x$ et $y$ sont \emph{disjoints} », qui
entraîne $x \neq y$ ;
\item[(4)] une relation symétrique plus fine, « $x$ et $y$ sont
\emph{adjacents} » ;
\item[(5)] une relation plus fine encore, non symétrique, « $x$ est
\emph{spécialisation} de $y$ », qui entraîne que $x$ est critique et $y$ ne
l'est pas ;
\item[(6)--(8)] un ensemble $S$ de \emph{segments}, et pour chaque $I \in S$
une paire $\delta I$ de lieux adjacents, ses extrémités, et une partie
$\widetilde{I} \supset \delta I$, les lieux sur le segment\footnote{Les
numéros (6) et (7) se chevauchent d'ajouts interlinéaires barrés ; l'ordre de
lecture est incertain, mais le contenu des données ne l'est pas. Le bord est
noté $\delta I$ sur ces deux pages, $\partial I$ ensuite.}.
\end{enumerate}

Trois remarques suivent, qui préparent les deux moutures suivantes. D'abord,
l'adjacence est redondante : $x$ et $y$ sont adjacents si et seulement s'il
existe un segment d'extrémités $\{x,y\}$. Ensuite --- « je présume » --- deux
lieux adjacents sont les extrémités d'au plus deux segments, et de deux
exactement lorsqu'ils sont sur un « anneau », une composante qui est un
cercle : c'est la situation que les pages 25 et 26 essaieront, à la fin du
dossier, de déduire des axiomes. Enfin, la question de savoir si un segment
est déterminé par $\delta I$ et $\widetilde{I}$ reste ouverte ; elle sera
tranchée à la page 3 (en marge, le 7 juin) et démontrée à la page 6.

Le seul axiome écrit, E$_{1}$, demande que sur $\mathcal{L}_{0}$ la relation
« $x$ et $y$ ne sont pas disjoints » soit une relation
d'équivalence\footnote{Entre « sont » et « disjoints », un signe barré qui
ressemble à $\neq$ ; la lecture « ne sont pas disjoints » est celle qui donne
un sens à la phrase suivante.} ; l'ensemble quotient est l'ensemble des
\emph{points} de la forme, et plusieurs lieux critiques peuvent définir le
même point critique. L'idée --- un point singulier vu comme une classe de
lieux qui « se touchent » --- est abandonnée par les moutures suivantes, qui
ne gardent comme données que les lieux et les segments ; elle reparaît en une
ligne à la page 9. Les pages 1 et 2 sont barrées chacune d'un long trait.

\pagerange{3}{4}
\section*{Deuxième mouture (pages 3 et 4)}

Les données se réduisent à deux ensembles $\mathcal{L}$ et $S$, une
application $S \to \mathfrak{P}(\mathcal{L})$, $I \mapsto \widetilde{I}$, et
une application $\partial : S \to \mathfrak{P}_{2}(\mathcal{L})$ avec
$\partial I \subset \widetilde{I}$. Une note datée du 7 juin décide, en
marge, d'identifier $I$ à $\widetilde{I}$ et $S$ à une partie de
$\mathfrak{P}(\mathcal{L})$\footnote{Le « $\widetilde{I}$ » de cette note est
lu avec doute. La date montre que la note est postérieure au corollaire 4 de
la page 6, qui la justifie.}.

Un lieu est \emph{régulier} s'il est intérieur à au moins un segment,
\emph{singulier} (critique, de branchement) sinon :
\[
\mathcal{L}_{r} = \bigcup_{I \in S} I^{\circ}, \qquad \mathcal{L}_{s} =
\mathcal{L} \smallsetminus \mathcal{L}_{r}, \qquad \widetilde{I} \cap
\mathcal{L}_{s} \subset \partial I .
\]
Un lieu singulier $y$ est une \emph{spécialisation} de $x$ s'il existe un
segment $I$ avec $x \in I^{\circ}$ et $y \in \partial I$ ; la donnée (5) de la
première mouture devient ainsi une notion définie, et le lieu critique reste
celui qui se spécialise.

Viennent des axiomes, sur une page presque entièrement annulée par une
grande croix\footnote{On résume ce qui se lit sous les traits. Deux axiomes
portent le numéro F$_{2}$, le second sur un numéro barré.} :
\emph{décomposition des segments} (tout lieu intérieur $z$ d'un segment
d'extrémités $x$, $y$ le coupe en deux sous-segments uniques $I_{x,z}$,
$I_{y,z}$ qui ne se rencontrent qu'en $z$) ; un axiome qui place tout lieu
intérieur dans l'intérieur d'un segment contenu dans $I^{\circ}$, d'où l'on
tire que $\widetilde{I}$ est infini et que $S$ l'est s'il n'est pas vide ;
\emph{recollement des segments}, réciproque de la décomposition, en un lieu
régulier ; enfin, toute partie finie de $I^{\circ}$ est contenue dans un
segment inclus dans $I^{\circ}$. Ce sont, sous d'autres numéros, les F1--F2,
F4 et F$'$5 de la troisième mouture. Deux notes de marge méritent d'être
retenues. L'une avertit : « Mais attention, il [\dots] dire $\widetilde{I} =
\widetilde{I}_{x,z} \cup \widetilde{I}_{y,z}$ !!! »\footnote{Deux mots sont
illisibles, et c'est d'eux que dépend le sens. La suite du dossier (voir la
section sur le recollement, pages 7 et 8) montre que l'égalité ne découle pas
des axiomes et est fausse dans les modèles qu'on a en vue ; nous lisons donc
la note comme une mise en garde contre cette égalité. C'est notre lecture.}.
L'autre --- lue « Apparemment » --- énonce que tout sous-segment de $I$ ayant
$z$ pour extrémité est contenu dans $\widetilde{I}_{x,z}$ ou dans
$\widetilde{I}_{y,z}$ : ce sera la seconde moitié de F2.

\pagerange{5}{6}
\section*{Troisième mouture : données et divisibilité (pages 5 et 6)}

Sous le titre « Modèle d'une forme $1$-dimensionnelle », les données sont
celles de la deuxième mouture : $\mathcal{L}$, $S$, $I \mapsto \widetilde{I}$,
$I \mapsto \partial I \in \mathfrak{P}_{2}(\mathcal{L})$, avec $\partial I
\subset \widetilde{I}$.

\begin{quote}
\textbf{F1} (axiome de divisibilité\footnote{Le mot « divisibilité », en
marge, est lu avec doute ; c'est celui que la page 12 emploie pour une
propriété voisine.}). \emph{Soient $I \in S$ et $\partial I = \{x,y\}$. Pour
tout $z \in \widetilde{I} \smallsetminus \{x\}$ il existe un unique $J \in S$,
noté $I_{x,z}$, tel que $\widetilde{J} \subset \widetilde{I}$ et $\partial J =
\{x,z\}$.}

\textbf{F2}. \emph{Pour tout $I \in S$, $I^{\circ} \neq \emptyset$ ; si
$\partial I = \{x,y\}$ et $z \in I^{\circ}$, alors $\widetilde{I}_{x,z} \cap \widetilde{I}_{y,z} = \{z\}$, et
tout $J \in S$ tel que $\widetilde{J} \subset \widetilde{I}$ et $z \in
\partial J$ vérifie $\widetilde{J} \subset \widetilde{I}_{x,z}$ ou
$\widetilde{J} \subset \widetilde{I}_{y,z}$.}
\end{quote}

L'unicité dans F1 donne $I_{x,y} = I_{y,x} = I$. F1 dit, de façon équivalente
(c'est la forme entre crochets de la page) : l'application qui envoie un
sous-segment $J$ de $I$ issu de $x$ sur son autre extrémité est une bijection
sur $\widetilde{I} \smallsetminus \{x\}$.

\begin{quote}
\textbf{Corollaires} (pages 5 et 6). \emph{Soit $I \in S$, $\partial I =
\{x,y\}$, et $S_{I} = \{J \in S \mid \widetilde{J} \subset \widetilde{I}\}$.}
\end{quote}

\begin{enumerate}
\item[1.] \emph{$I^{\circ}$ est infini.}
\item[2.] \emph{Dans F2, les deux inclusions s'excluent.}
\item[3.] \emph{L'application $S_{I} \to \mathfrak{P}_{2}(\widetilde{I})$, $J
\mapsto \partial J$, est injective ; son image contient toutes les paires qui
rencontrent $\partial I$ ; une paire $\{u,v\} \subset I^{\circ}$ est dans
l'image si et seulement si $v \in \widetilde{I}_{x,u} \cup
\widetilde{I}_{y,u}$, si et seulement si $u \in \widetilde{I}_{x,v} \cup
\widetilde{I}_{y,v}$.}
\item[4.] \emph{Si $\widetilde{I} = \widetilde{J}$, alors $I = J$.}
\end{enumerate}

La page ne démontre que le premier, d'un mot (« découle de la condition de
divisibilité »). Les quatre sont justes\footnote{Démonstrations, qui ne sont
pas sur la page. (1) Si $z \in I^{\circ}$, F2 appliqué à $I_{x,z}$ donne $z'
\in I_{x,z}^{\circ}$ ; l'unicité dans F1 montre que $I_{x,z'}$ est le
sous-segment de $I_{x,z}$ d'extrémités $x$, $z'$, donc
$\widetilde{I}_{x,z'} \subsetneq \widetilde{I}_{x,z}$, et $z \notin
\widetilde{I}_{x,z'}$ par F2 dans $I_{x,z}$ : on construit ainsi une suite
infinie de lieux distincts de $I^{\circ}$. (2) $\widetilde{J}$ contient les
deux éléments de $\partial J$, et $\widetilde{I}_{x,z} \cap
\widetilde{I}_{y,z} = \{z\}$. (3) Si $\partial J = \partial J' = \{u,v\}$ avec
$u \in I^{\circ}$, F2 place $J$ et $J'$ dans le même $I_{x,u}$ ou $I_{y,u}$
(sinon $v$ serait dans les deux, donc égal à $u$), et l'unicité de F1 dans ce
sous-segment donne $J = J'$ ; si $u \in \partial I$, c'est F1. Les
équivalences de c) suivent de F2 et de F1. (4) Si $\partial J$ rencontre
$\partial I$, en $x$ disons, $J = I_{x,z}$ ; si $z \neq y$, F2 donne $y \notin
\widetilde{I}_{x,z} = \widetilde{J} = \widetilde{I}$, absurde. Si $\partial J
\subset I^{\circ}$, F2 place $\widetilde{J}$ dans $\widetilde{I}_{x,u}$ ou
$\widetilde{I}_{y,u}$, qui ne contiennent pas tout $\widetilde{I}$.}.

Le point c) mérite qu'on s'y arrête, parce qu'il dit ce que le modèle
\emph{n'est pas}. Deux lieux intérieurs $u$, $v$ d'un segment ne sont pas
toujours les extrémités d'un sous-segment : il faut que l'un soit « du côté
de $x$ » ou « du côté de $y$ » de l'autre. La page ajoute : « on sait donc que
$u$ et $v$ sont des lieux \emph{disjoints} sur le segment $I$ », et promet
d'y revenir\footnote{« y reviendra » est lu avec doute ; le « ou » de la
condition 2), souligné deux fois, aussi.}. C'est le sens que le mot
« disjoints » a dans le dossier 156-2 (page 8) : deux lieux d'une faille sont
disjoints quand ils sont emboîtés, l'un du côté de $A$ par rapport à l'autre.
Rien n'oblige deux lieux d'un même segment à l'être.

Le corollaire 4 autorise ce que la marge de la page 3 avait décidé :
\emph{désormais on identifie un segment à l'ensemble de ses lieux, et l'on
écrit $I$ pour $\widetilde{I}$}, $I^{\circ} = I \smallsetminus \partial I$.

\pagerange{6}{7}
\section*{Branches (pages 6 et 7)}

\begin{quote}
\textbf{F3} (axiome des branches). \emph{Soient $I, J \in S$ et $x \in
\partial I \cap \partial J$. Il existe $u \in I \smallsetminus \{x\}$ et $v
\in J \smallsetminus \{x\}$ tels que, ou bien $I_{x,u} = J_{x,v}$ (et alors
$u = v$), ou bien $I_{x,u} \cap J_{x,v} = \{x\}$.}
\end{quote}

Les deux cas s'excluent. La page écrit d'abord $u \in I^{\circ}$, $v \in
J^{\circ}$, et la marge corrige en $u \in I \smallsetminus \{x\}$, $v \in J
\smallsetminus \{x\}$\footnote{« devrait plutôt [être] $u \in I
\smallsetminus \{x\}$ » ; la fin de la note, qui renvoie à la divisibilité, est
illisible. Lorsque $I^{\circ}$ et $J^{\circ}$ ne sont pas vides, les deux
formes sont équivalentes (on rétrécit par F1) ; la forme de la marge est la
seule qui ait un sens pour les réseaux réduits de la page 23, et c'est elle
que nous retenons.}. Un segment issu de $x$ et un autre, ou bien commencent
par un même morceau, ou bien partent de $x$ dans des directions qui ne se
recouvrent pas près de $x$.

Soit $S_{x}$ l'ensemble des segments dont $x$ est une extrémité. On dit que
$I, J \in S_{x}$ définissent la \emph{même branche} en $x$, et l'on écrit $I
\mathrel{R_{x}} J$, s'il existe $z \in I^{\circ} \cap J^{\circ}$ avec $I_{x,z}
= J_{x,z}$. La page affirme que c'est une relation d'équivalence ; elle est
réflexive et symétrique, mais sa transitivité demande une condition que F1 à
F3 ne donnent pas :
\begin{quote}
\textbf{F$'$3} (filtration en une extrémité). \emph{Si $I \in S$, $x \in
\partial I$ et $u, v \in I^{\circ}$, il existe $w \in I^{\circ}$ avec $I_{x,w}
\subset I_{x,u} \cap I_{x,v}$.}
\end{quote}
Une note serrée au bas de la page 6 commence justement : « Il manque [\dots]
à partir de (F$'_{3}$) : Soit $I \in S$, $u, v \in I^{\circ}$, alors [\dots]
$\exists\, w$ [\dots] $I_{x,w}$ »\footnote{La note est en grande partie
illisible ; l'énoncé F$'$3 ci-dessus est notre reconstruction, guidée par ce
qui s'en lit et par ce dont la transitivité a besoin. Avec F$'$3 : si $I
\mathrel{R_{x}} J$ par $z$ et $J \mathrel{R_{x}} K$ par $z'$, on prend $w$
dans $J^{\circ}$ avec $J_{x,w} \subset J_{x,z} \cap J_{x,z'}$ ; l'unicité de
F1 donne $I_{x,w} = J_{x,w} = K_{x,w}$, et $w$ est intérieur à $I$ et à $K$.}.
Avec F$'$3, $R_{x}$ est une équivalence. L'ensemble quotient
\[
B_{x} = S_{x} / R_{x}
\]
est l'ensemble des \emph{branches} en $x$, et l'on note $\mathrm{br}_{x}(I)$
la branche de $I$. Le cardinal de $B_{x}$ est l'\emph{ordre} de $x$ : un lieu
d'ordre $0$ est \emph{isolé}, un lieu d'ordre $1$ est un \emph{bout}. C'est,
pour un graphe, le degré d'un sommet\footnote{Le rapprochement est le nôtre ;
la page 23 parlera d'« ordre ou multiplicité d'un nœud ». La page écrit
ensuite « Ordre fini », seul sur sa ligne : peut-être une exigence laissée en
suspens ; on ne l'impose pas.}.

\pagerange{7}{8}
\section*{Lieux réguliers et recollement (pages 7 et 8)}

Les lieux réguliers, singuliers, sont définis comme à la page 3.

\begin{quote}
\textbf{F4} (axiome de recollement). \emph{Soient $I', I'' \in S$ et $z$ un
lieu \emph{régulier} tels que $I' \cap I'' = \partial I' \cap \partial I'' =
\{z\}$ ; posons $\partial I' = \{x,z\}$, $\partial I'' = \{y,z\}$. Il existe
un unique $I \in S$ tel que $I \supset I' \cup I''$ et $\partial I =
\{x,y\}$ ; alors $I' = I_{x,z}$ et $I'' = I_{y,z}$.}
\end{quote}

La page ajoute entre parenthèses une remarque sur l'égalité, dont un mot
manque\footnote{« (on \uncertain{a} [\dots] l'égalité !) », le verbe lu avec
doute et le mot suivant illisible. Lu « on a l'égalité », ce serait faux en
général, et nous ne l'affirmons pas.}. Il ne faut pas lire $I = I' \cup I''$ :
on verra aux pages 10 à 13 qu'un segment est un ensemble ordonné, et que
$I_{x,z} \cup I_{z,y}$ est l'ensemble des lieux de $I$ comparables à $z$ ;
l'égalité revient donc à demander que l'ordre soit total, ce que ni les
axiomes ni les modèles du chapitre II n'imposent. Seule l'inclusion est vraie,
et c'est elle qu'on énonce.

\begin{quote}
\textbf{Corollaire 1} (page 7). \emph{Soient $I \in S$, $\partial I =
\{x,y\}$, avec $x$ régulier. Il existe $J \in S$ tel que $J \supset I$,
$\partial J \cap \partial I = \{y\}$ et $x \in J^{\circ}$.}
\end{quote}

Le segment se prolonge au-delà d'une extrémité régulière. L'esquisse de la
page : F1 et F2 montrent que $x$ est d'ordre $\geq 2$ ; F3 fournit alors un
segment $K$ issu de $x$ avec $K \cap I = \partial K \cap \partial I = \{x\}$ ;
F4 recolle $K$ et $I$\footnote{Le premier pas est juste : si $x \in
L^{\circ}$, les sous-segments $L_{a,x}$ et $L_{b,x}$ ne se rencontrent qu'en
$x$, donc ne définissent pas la même branche. Le deuxième ne l'est pas tel
quel : appliqué à $I$ et à un segment d'une autre branche, F3 donne un $K$ qui
ne rencontre que $I_{x,u}$ en $\{x\}$, non $I$ tout entier ; un segment issu
de $x$ peut revenir rencontrer $I$ loin de $x$, comme sur un cercle. Il faut
pouvoir rétrécir $K$ jusqu'à éviter $I \smallsetminus \{x\}$ ; c'est vrai dans
les modèles qu'on a en vue, mais cela ne découle pas visiblement des axiomes
et la page ne le fait pas. La fin de la ligne, « avec $y' \neq x$ », est lue
avec doute.}.

\begin{quote}
\textbf{Corollaire 2} (page 8). \emph{Si les deux extrémités de $I$ sont
régulières, il existe $J \in S$ tel que $I \subset J^{\circ}$.}
\end{quote}

C'est le corollaire 1 appliqué deux fois, et il en hérite la réserve.

\begin{quote}
\textbf{F5} (axiome des lieux réguliers). \emph{Un lieu régulier est d'ordre
$2$.}

\textbf{F$'$5}. \emph{Pour $x$ régulier, les $I^{\circ}$ qui contiennent $x$
forment une base de filtre : si $x \in I'^{\circ} \cap I''^{\circ}$, il
existe $I$ tel que $x \in I^{\circ} \subset I'^{\circ} \cap I''^{\circ}$.}
\end{quote}

On sait déjà que l'ordre d'un lieu régulier est au moins $2$ ; F5 interdit le
branchement en un lieu régulier, ce qui justifie a posteriori que les
lieux singuliers s'appellent aussi « de branchement ». La page affirme que F5
et F$'$5 sont équivalents, sans démonstration\footnote{Nous ne reproduisons
pas l'équivalence comme acquise. Une note marginale de six ou sept lignes, en
face des deux énoncés, est presque entièrement illisible ; on y lit F$'_{5}$,
F$_{2}$, F$_{5}$ et « $[0,1]$ ». Le NB qui suit F$'$5 (« $\Longrightarrow I
\subset I'^{\circ} \cap I''^{\circ}$ ») se lit au mieux : quitte à rétrécir
$I$, on peut supposer le segment fermé $I$ lui-même contenu dans
$I'^{\circ} \cap I''^{\circ}$.}. Dans un modèle, F$'$5 dit que les
voisinages de $x$ formés d'intérieurs de segments sont filtrants, comme les
intervalles ouverts autour d'un point de la droite.

\pagerange{8}{9}
\section*{Restriction à une partie (pages 8 et 9)}

Une première remarque (page 8), barrée, demande si les axiomes passent à la
réunion $\mathcal{L}'$ des lieux d'une famille de segments. La page 9 reprend
plus généralement :

\emph{Remarque 1.} Les axiomes sont trivialement satisfaits si $S =
\emptyset$, c'est-à-dire si tous les lieux sont singuliers ; un modèle se
réduit alors à un ensemble. La page annonce qu'on enrichira la structure en
introduisant une relation d'équivalence sur $\mathcal{L}_{s}$ --- l'idée de
l'axiome E$_{1}$ de la page 2, un point singulier comme classe de lieux. Elle
n'y revient pas.

\emph{Remarque 2.} Pour une partie quelconque $\mathcal{L}' \subset
\mathcal{L}$, posons $S' = \{I \in S \mid I \subset \mathcal{L}'\}$. La page
affirme que $(\mathcal{L}', S')$ satisfait F1 à F5. C'est juste pour F1, F2,
F3 (et F$'$3), F5 et F$'$5 ; pour F4, il faut que le segment recollé reste
dans $\mathcal{L}'$\footnote{Démonstration, qui n'est pas sur la page : les
sous-segments d'un segment de $S'$ sont dans $S'$, ce qui règle F1 à F3 ; une
branche en $x$ pour $S'$ est contenue dans une branche pour $S$, d'où un ordre
$\leq 2$ en un lieu régulier de $\mathcal{L}'$, et $\geq 2$ par l'argument du
corollaire 1 ; le segment $I$ fourni par F$'$5 est contenu dans $I'$. Pour
F4, le recollé $I$ de $I', I'' \in S'$ contient $I' \cup I''$ mais peut
contenir d'autres lieux, extérieurs à $\mathcal{L}'$ : l'hérédité de F4 est
vraie si $I = I' \cup I''$, c'est-à-dire, d'après la section précédente,
quand les ordres sont totaux. C'est une hypothèse que la page ne fait pas.}.
Le cas qui servira est celui où $\mathcal{L}' = I$ est un segment : alors $S'
= S_{I}$, qu'on identifie par le corollaire 3 à une partie de
$\mathfrak{P}_{2}(I)$\footnote{La page écrit « les segments tels que
$\widetilde{I} \subset I$ » pour $\widetilde{J} \subset \widetilde{I}$. Dans
ce cas F4 est hérité : on le vérifie avec l'ordre de la section suivante, où
le recollé de $[x,z]$ et $[z,y]$ est $[x,y] \subset I$.}.

\pagerange{10}{14}
\section*{Le modèle-segment est un biordre (pages 10 à 14)}

Un modèle est un \emph{modèle-segment} s'il existe $I \in S$ avec $I =
\mathcal{L}$. Ce segment est unique (corollaire 4), et $\mathcal{L}_{r} =
I^{\circ}$. Choisissons une extrémité $x_{0}$ et notons $x_{1}$ l'autre. Une
première tentative de définir un ordre (page 10) est cernée et
barrée\footnote{Elle définit $x \preccurlyeq y$ par $y \in [x,x_{1}]$, prouve
la transitivité (« c'est immédiat ») et commence une proposition qui
s'interrompt. La page 11 la reprend avec une meilleure présentation.
« mutuellement » est lu avec doute.} ; la page 11 pose
\[
[x_{0},x] = \begin{cases} \{x_{0}\} & x = x_{0}, \\ I_{x_{0},x} & x \neq x_{0},
\end{cases}
\]
et de même $[x,x_{1}]$ en échangeant $x_{0}$ et $x_{1}$.

\begin{quote}
\textbf{Proposition} (page 11). \emph{Pour $x, y \in \mathcal{L}$, les
conditions suivantes sont équivalentes : (i) $[x_{0},x] \subset [x_{0},y]$ ;
(ii) $x \in [x_{0},y]$ ; (i bis) $[x,x_{1}] \supset [y,x_{1}]$ ; (ii bis) $y
\in [x,x_{1}]$. La relation $x \preccurlyeq y$ ainsi définie est un ordre sur
$\mathcal{L}$.}
\end{quote}

C'est un préordre, évidemment sur (i) ; c'est un ordre parce que $[x_{0},x] =
[x_{0},y]$ entraîne, par le corollaire 4 et en passant aux extrémités, $x =
y$\footnote{L'équivalence de (ii) et (ii bis), qui n'est pas sur la page, vient
du corollaire 3 c) et de F2 : si $x \in I_{x_{0},y}$, il y a un sous-segment
d'extrémités $x$, $y$ ; F2 en $x$ le place dans $I_{x_{0},x}$ ou dans
$I_{x_{1},x}$, et le premier cas donnerait $I_{x_{0},x} = I_{x_{0},y}$, donc $x
= y$.}. Désormais on écrit $\leq$ et $<$, comme la page à partir de la page 12.

\begin{quote}
\textbf{Proposition} (pages 12 et 13). \emph{L'ordre $\leq$ a un plus petit
élément $x_{0}$ et un plus grand élément $x_{1}$, distincts. Les segments du
modèle sont exactement les $[x,y] = \{z \mid x \leq z \leq y\}$ pour $x < y$,
avec $\partial [x,y] = \{x,y\}$. L'ordre est}
\end{quote}

\begin{enumerate}
\item[a)] \emph{dense : si $x < y$, il existe $z$ avec $x < z < y$ ;}
\item[b)] \emph{localement filtrant décroissant : si $x < y_{1}$ et $x <
y_{2}$, il existe $y$ avec $x < y \leq y_{1}, y_{2}$ ;}
\item[b$'$)] \emph{localement filtrant croissant : si $y_{1} < x$ et $y_{2} <
x$, il existe $y$ avec $y_{1}, y_{2} \leq y < x$.}
\end{enumerate}

\begin{quote}
\emph{Réciproquement, un ensemble ordonné ayant un plus petit et un plus grand
élément distincts, et vérifiant a), b), b$'$), définit, avec ces segments, un
modèle-segment.}
\end{quote}

La page énonce ces propriétés et la réciproque sans
démonstration\footnote{L'énumération de la page 12 est surchargée
d'ajouts ; « filtrante croissante » et l'astérisque qui renvoie à la marge
sont lus en partie. Esquisse : que les segments soient les $[x,y]$ vient de
F2, du corollaire 3 et de F4 (si $x < z < y$, les segments $[x,z]$ et $[z,y]$
ne se rencontrent qu'en $z$, régulier, et leur recollé est le segment
d'extrémités $x$, $y$) ; a) est F2 ; b) est F$'$3 appliqué au segment
$[x,x_{1}]$ en son extrémité $x$, et b$'$) de même avec $[x_{0},x]$. Pour la réciproque, F1 à
F5 se vérifient directement sur les intervalles, F3 et F$'$3 grâce à b) et
b$'$).}. Ce qu'elle n'énonce pas, et qu'il faut souligner, c'est que l'ordre
soit total : il ne l'est pas en général, et c'est le point c) de la page 5 vu
autrement. Le cardinal de $\mathcal{L}$ est infini, par la densité.

Échanger $x_{0}$ et $x_{1}$ remplace l'ordre par l'ordre opposé, et les
conditions a), b), b$'$) sont autoduales. D'où la conclusion de la page 13 :
\emph{une structure de modèle-segment équivaut à une structure de biordre} ---
une paire d'ordres opposés --- dont chacun a un plus petit et un plus grand
élément distincts, et qui est dense et localement filtrant dans les deux sens.
Une \emph{orientation} du segment (page 14) est le choix d'une extrémité,
c'est-à-dire de l'un des deux ordres.

Ce sont mot pour mot les propriétés a) à c) que la page 9 du dossier 156-2,
écrite le lendemain, énonce pour l'ensemble des lieux d'une faille
topologique orientée : le chapitre II fournit au chapitre I son modèle
géométrique\footnote{Le rapprochement est le nôtre ; aucun des deux dossiers
ne renvoie à l'autre sur ce point. Le cas où les lieux sont des points, donc
où l'ordre est total, est celui du segment $[0,1]$ ; mais sur une bande
$S^{1} \times [0,1]$, deux cercles ondulés qui se croisent sont deux lieux
non comparables, et l'ensemble des lieux n'est pas totalement ordonné. La
caractérisation d'un intervalle par des axiomes d'ordre ou
d'« intermédiarité » (\emph{betweenness}) remonte à M.~Pasch (1882),
E.~V.~Huntington et D.~Hilbert ; un ordre dense non nécessairement total,
localement filtrant (en anglais \emph{locally directed}), n'a pas de nom
consacré.}.

\pagerange{14}{17}
\section*{Modèles réguliers connexes, orientation (pages 14 à 17)}

Un modèle est \emph{régulier} si $\mathcal{L} = \mathcal{L}_{r}$. Sur
$\mathcal{L}$, la relation « il existe un segment contenant $x$ et $y$ » est
réflexive et symétrique ; ses classes pour l'équivalence engendrée sont les
\emph{composantes connexes}, et leur ensemble est noté
$\pi_{0}(\mathfrak{X})$\footnote{La page invoque le corollaire 2 pour dire
que, dans un modèle régulier, on peut demander $x, y \in I^{\circ}$ ; c'est
vrai sous la réserve faite sur le corollaire 1. La définition de
$\pi_{0}$ n'en dépend pas.}. Tout modèle est la somme disjointe de ses
composantes, et les composantes d'un modèle régulier sont régulières : on se
ramène aux modèles réguliers connexes.

\begin{quote}
\textbf{Théorème} (énoncé pages 15 à 17). \emph{Soit $\mathfrak{X}$ un modèle
régulier connexe. Il existe, de façon unique, un ensemble $\omega$ à deux
éléments et des bijections $\omega \simeq \partial I$ ($I \in S$) et $\omega
\simeq B_{x}$ ($x \in \mathcal{L}$) telles que :}
\end{quote}

\begin{enumerate}
\item[a)] \emph{si $J \subset I$ sont deux segments, la bijection $\partial I
\simeq \partial J$ qui en résulte envoie le plus petit élément de $I$ sur le
plus petit élément de $J$, pour l'un quelconque des deux ordres de $I$ (et
l'ordre qu'il induit sur $J$) ;}
\item[b)] \emph{si $x$ est le plus petit élément de $I$ pour l'un de ses
ordres, l'élément $x \in \partial I$ correspond à la branche
$\mathrm{br}_{x}(I)$.}
\end{enumerate}

Autrement dit : sur la réunion disjointe
\[
\coprod_{I \in S} \partial I \;\amalg\; \coprod_{x \in \mathcal{L}} B_{x},
\]
soit $R$ la relation d'équivalence engendrée par $(I,x) \sim (J,y)$ lorsque
$J \subset I$ et que $x$, $y$ se correspondent comme en a), et par $(I,x) \sim
(x, \mathrm{br}_{x}(I))$. Le théorème affirme que $R$ a exactement deux
classes, et que les deux extrémités d'un même segment tombent dans des classes
différentes. La page ne le démontre pas : « [\dots] vérifier ici. Si c'était
faux, il faudrait le mettre en axiome. S'il est vrai, il doit être immédiat
que le nombre de classes est $1$ ou $2$. »\footnote{Le début de la phrase est
illisible, et l'ajout de deux lignes en tête de la page 17 n'est lu qu'en
partie. La page écrit $\omega_{\mathfrak{X}}$ puis $\omega_{\mathcal{L}}$ pour
le même quotient ; nous écrivons $\omega$. Que le nombre de classes soit au
plus $2$ découle de la connexité ; qu'il soit exactement $2$ --- que le modèle
soit orientable --- est la partie qui reste à établir ou à poser en axiome. Le
croquis de marge montre un arc portant un point épais et deux flèches dans le
même sens.}

L'ensemble $\omega$ est l'ensemble des \emph{orientations} de $\mathfrak{X}$.
En termes actuels, les bijections $\omega \simeq \partial I$, $\omega \simeq
B_{x}$ trivialisent le torseur des orientations, et le théorème dit que le
revêtement d'orientation est trivial\footnote{Le vocabulaire est le nôtre.
Pour une variété de dimension $1$ connexe --- une droite ou un cercle --- le
revêtement d'orientation à deux feuillets est toujours trivial ; c'est ce que
le théorème transporte dans le cadre combinatoire, où il n'est pas
évident.}. Une orientation choisie, chaque segment a une \emph{origine}
$\mathrm{or}(I)$ et une \emph{extrémité} $\mathrm{ex}(I)$, et chaque lieu une
branche sortante\footnote{« ainsi », « orienté » et « sortant » sont lus avec
doute.}. On pose
\[
x \prec y \iff \text{il existe } I \in S \text{ avec } \mathrm{or}(I) = x,\
\mathrm{ex}(I) = y,
\]
et l'on distingue le \emph{cas linéaire}, où $\prec$ est transitive, du
\emph{cas circulaire}, où elle ne l'est pas. Seul le premier est étudié.

\pagerange{17}{21}
\section*{Le cas linéaire, et la filtration relative (pages 17 à 21)}

Dans le cas linéaire, $x \preccurlyeq y \iff (x = y \text{ ou } x \prec y)$
est un préordre, et c'est un ordre : $x \prec y \prec x$ donnerait $x \prec
x$, impossible puisque les deux extrémités d'un segment sont distinctes. Les
segments sont les $[x,y]$ pour $x \prec y$, d'origine $x$ ; la structure de
modèle connexe orienté régulier est donc connue quand on connaît l'ordre.
La page affirme que cet ordre est
\begin{enumerate}
\item[--] \emph{strictement filtrant croissant} : pour tous $x, y$, il existe
$z$ avec $x, y < z$ ; et strictement filtrant décroissant ;
\item[--] localement strictement filtrant dans les deux sens : si $x < y$ et
$x < z$, il existe $u$ avec $x < u \leq y, z$, et dualement\footnote{La page
écrit « si $x \leq y, z$, $\exists\, u$ tel que $x \leq u \leq y, z$ », ce qui
est trivial ($u = x$) ; nous lisons la version stricte, qui est celle des
pages 12 et 20. Les formules de la parenthèse sont surchargées.} ;
\end{enumerate}
et note en marge qu'il n'a donc ni plus grand ni plus petit élément, et qu'il
est dense.

\emph{La démonstration qui bute} (pages 19 et 20). Pour montrer que l'ordre
est filtrant croissant, la page relie deux lieux $x$, $y$ par une chaîne
$x = x_{0}, x_{1}, \dots, x_{n} = y$ avec $x_{i}, x_{i+1} \in I_{i}^{\circ}$ ;
en posant $y_{i} = \mathrm{or}(I_{i})$ on a $x_{i} \geq y_{i} \leq
x_{i+1}$\footnote{Les signes d'inégalité sont repassés ; l'alternance est lue
avec doute, mais c'est la seule qui suive de la définition.}, de sorte que
$\mathcal{L}$ est connexe en tant qu'ensemble ordonné. On raisonne par
récurrence sur la longueur de la chaîne. Pour $n \leq 1$, il suffit qu'il y
ait un lieu au-dessus de $y$, ce qu'assure la régularité de $y$. Au pas de
récurrence, on a $x'$ au-dessus de $x$ et de $z$, et $y$ voisin de $z$ ; si
$z \geq y$ c'est fini, et si $z \leq y$, il faut trouver $t > x', y$ sachant
seulement que $z < x'$ et $z < y$. La page isole ce qui manque sous le nom de
\emph{filtration relative} :
\[
u < v \ \text{ et } \ u < w \quad \Longrightarrow \quad \text{il existe } t
\text{ avec } v, w < t,
\]
deux croquis montrant le « V » $v > u < w$ qu'on veut fermer en losange.

Ce n'est pas une conséquence des propriétés locales, et la page le dit
aussitôt : un ordre connexe, localement strictement filtrant dans les deux
sens et sans plus grand ni plus petit élément définit un modèle régulier
connexe --- on vérifie F1 à F5 --- mais n'est pas nécessairement filtrant
croissant\footnote{La phrase se termine sur un mot lu « (dites) », d'une
lecture incertaine.}. Le contre-exemple est à la page 21.

\begin{quote}
\textbf{Exemple} (page 21). \emph{Soient $\mathcal{L}_{0}$ et
$(\mathcal{L}_{i})_{i \in I}$, $\operatorname{card} I \geq 2$, des ensembles
ordonnés non vides, localement strictement filtrants dans les deux sens, et
tels que tout élément $a$ ait un $a' < a$ et un $a'' > a$. Sur $\mathcal{L} =
\mathcal{L}_{0} \amalg \coprod_{i} \mathcal{L}_{i}$, on pose $x < y$ si $x, y$
sont dans un même $\mathcal{L}_{i}$ ($i \in I$ ou $i = 0$) et $x < y$ dans
celui-ci, ou si $x \in \mathcal{L}_{0}$ et $y \in \mathcal{L}_{i}$, $i \in
I$.}
\end{quote}

C'est un arbre qui se ramifie vers le haut : un tronc $\mathcal{L}_{0}$
au-dessus duquel partent plusieurs rameaux. L'ordre satisfait les mêmes
conditions locales, mais deux lieux pris dans deux rameaux différents n'ont
aucun majorant : il n'est pas filtrant croissant\footnote{La page écrit
« $\exists\, x' < a < a''$ », que nous lisons « pour tout $a$, il existe $a' <
a < a''$ » ; « non vide » est inséré au-dessus de la ligne et lu avec doute.
Il faut $\operatorname{card} I \geq 2$, que la page ne dit pas, pour que
l'exemple en soit un ; il faut aussi $\mathcal{L}_{0}$ sans plus grand
élément, ce que donne $a''$. Un premier exemple, sur $\mathbb{Q}$ coupé en un
irrationnel $r$ avec une partie supérieure multipliée par un ensemble $E$ de
cardinal $\geq 2$, est barré de quatre traits et s'interrompt.}. Dans le
modèle correspondant, le tronc se dédouble sans qu'aucun lieu marque la
bifurcation : chaque lieu est régulier, d'ordre $2$, et pourtant la forme
bifurque. F1 à F5 ne voient pas cette bifurcation, parce qu'elle a lieu
« entre » des lieux, en un endroit où il n'y a pas de lieu.

\pagerange{22}{22}
\section*{L'axiome F6 (page 22)}

\begin{quote}
\textbf{F6}. \emph{Soient $I, J \in S$ et $x \in \partial I \cap \partial J$
définissant la même branche en $x$. Il existe $K \in S$ tel que $x \in
\partial K$ et $I^{\circ} \cup J^{\circ} \subset K^{\circ}$.}
\end{quote}

En d'autres termes, les intérieurs des segments issus de $x$ dans une branche
donnée forment une famille filtrante croissante pour l'inclusion ; la page
renvoie à F3, dont c'est le pendant « vers l'extérieur ». C'est exactement ce
qui exclut l'exemple de la page 21 : deux segments $[x,v]$ et $[x,w]$ issus
d'un lieu $x$ du tronc et montant dans deux rameaux définissent la même
branche en $x$, et aucun segment ne les contient tous deux.

La page en tire, pour $I, J \in S$ et $x \in \partial I \cap \partial J$ :
\begin{enumerate}
\item[a)] \emph{$\mathrm{br}_{x}(I) = \mathrm{br}_{x}(J)$ si et seulement s'il
existe $K \in S$ tel que $x \in \partial K$ et $I \cup J \subset K$ ;}
\item[b)] \emph{sinon, il existe $u \in I^{\circ}$, $v \in J^{\circ}$ tels que
$I_{x,u} \cap J_{x,v} = \{x\}$.}
\end{enumerate}

b) est F3\footnote{La page écrit $I_{x,u} \cap I_{x,v}$, pour $J_{x,v}$.}. Dans
a), la réciproque est juste (avec F$'$3) ; mais F6 ne donne $K$ qu'avec
$I^{\circ} \cup J^{\circ} \subset K^{\circ}$, et rien n'assure que les
extrémités autres que $x$ soient dans $K$\footnote{La note de marge, très
raturée, porte justement là-dessus : « On n'exige pas [\dots] que les
intersections $I \cap J$ [\dots] soient dans $K$. Mais [\dots] celles de ces
extrémités [qui sont] régulières [\dots] (cor.~2 p.~7) » --- c'est-à-dire,
dans notre pagination, le corollaire 2 de la page 8. On y lit l'idée
d'utiliser ce corollaire pour faire entrer dans $K$ les extrémités régulières.
« On a donc », qui introduit a) et b), est lu avec doute.}. C'est cette forme
fermée qu'il faut pour achever la filtration relative de la page 20 (avec $K
= [u,k]$, on veut $v, w \leq k$) et pour l'argument de la page 25. Nous
retenons donc F6 tel qu'il est écrit, et a) sous la réserve que la forme
fermée soit démontrée par la voie qu'indique la marge, ou prise comme axiome.

\pagerange{23}{24}
\section*{Réseaux (pages 23 et 24)}

Le 7 juin, l'objet reçoit son nom : un modèle de $1$-forme satisfaisant F1 à
F6 est un \emph{réseau}. Ses lieux singuliers sont ses \emph{nœuds} ; les
composantes connexes de $\mathcal{L}_{r}$ sont ses \emph{lignes} --- « de
préférence : arêtes »\footnote{La page écrit « $\mathcal{L}_{r} \smallsetminus
\mathcal{L}_{s}$ », qui est $\mathcal{L}_{r}$ puisque les deux ensembles sont
disjoints ; l'indice $r$ est d'ailleurs lu avec doute. On entend les
composantes du sous-modèle $(\mathcal{L}_{r}, S')$ de la page 9.} ; ses
composantes connexes sont ses \emph{composantes}. L'ordre, ou
\emph{multiplicité}, d'un nœud est le nombre de ses branches, égal à $2$ en
un lieu régulier.

Un réseau est \emph{divisible} si $I^{\circ} \neq \emptyset$ pour tout
segment, \emph{réduit} si $I^{\circ} = \emptyset$ pour tout segment,
c'est-à-dire $I = \partial I$, ou encore $\mathcal{L}_{r} = \emptyset$ : tous
ses lieux sont des nœuds. Ces deux définitions supposent tacitement qu'on a
retiré de F2 sa première clause, $I^{\circ} \neq \emptyset$, sans quoi tout
réseau serait divisible et aucun ne serait réduit (à moins que $S$ soit
vide)\footnote{La page ne le dit pas. Nous lisons donc « réseau » avec F2
privé de cette clause, et F3 sous la forme de la marge de la page 6 ;
« divisible » restitue la clause. Le mot qui suit « divisible » est
illisible.}.

\begin{quote}
\emph{La donnée d'une structure de réseau réduit sur $\mathcal{L}$ revient à
celle d'une partie $S$ de $\mathfrak{P}_{2}(\mathcal{L})$} --- « les axiomes
F1 à F6 sont vides ! » --- \emph{donc à une structure de graphe réduit.}
\end{quote}

C'est juste sous la lecture qu'on vient de dire\footnote{Avec $I^{\circ} =
\emptyset$, F2, F4, F5 et F6 ne portent que sur des lieux intérieurs ou
réguliers, et sont vides ; F1 se réduit à $I_{x,y} = I$ ; F3, sous la forme de
la marge, est satisfait avec $u$, $v$ les autres extrémités, car deux
segments distincts issus de $x$ n'ont que $x$ en commun. C'est l'identification
de $S$ à une partie de $\mathfrak{P}(\mathcal{L})$ qui exclut les arêtes
multiples.}. En termes actuels, un graphe réduit est un \emph{graphe simple}
: pas de boucle, puisqu'un segment a deux extrémités distinctes, pas d'arête
multiple, puisqu'un segment est déterminé par ses lieux. Le cercle à un ou
deux sommets, qui demanderait une boucle ou deux arêtes parallèles, n'est donc
pas un réseau réduit ; il en faut au moins trois nœuds. Le lien avec les
graphes est le but visible de la notion : un réseau général est un graphe
dont les arêtes sont des segments divisibles.

\pagerange{24}{26}
\section*{Deux segments consécutifs (pages 24 à 26)}

« Je reviens », dit la page, à l'étude des réseaux réguliers, connexes et non
vides. Soient deux segments consécutifs $I$, d'extrémités $x$, $y$, et $J$,
d'extrémités $y$, $z$, avec $y$ régulier et $\mathrm{br}_{y}(I) \neq
\mathrm{br}_{y}(J)$\footnote{Un « Lemme » est commencé puis barré. La
première lettre de la note de marge « $\in \partial I \cap \partial J$ » peut
se lire $x$ ou $y$ ; la figure suggère $y$.}.

\emph{Cas I : $I \cap J = \{y\}$.} Alors $x \neq z$, et l'axiome de
recollement donne un unique $K \supset I \cup J$, d'extrémités $x$, $z$, avec
$I = K_{x,y}$, $J = K_{y,z}$. « Situation bien comprise. »

\emph{Cas II : $I \cap J \neq \{y\}$.} Un premier traitement est barré (page
24) et repris (page 25). Si $w \in I \cap J$, $w \neq y$, les sous-segments
$I_{w,y}$ et $J_{y,w}$ ont les mêmes extrémités $\{w,y\}$ et des branches
différentes en $y$\footnote{La marge note que $w$ est régulier, sauf si $w =
x = z$. Les croquis de ces deux pages sont des cercles : deux arcs de mêmes
extrémités. Les figures barrées ne sont pas reprises.}. La page étudie donc
d'abord la situation où $\partial I = \partial J = \{x,y\}$ et
$\mathrm{br}_{y}(I) \neq \mathrm{br}_{y}(J)$, sans plus supposer $y$ régulier,
et demande ce qu'est $I \cap J$.

\emph{Les branches diffèrent aussi en $x$.} Sinon, F6 fournirait $K$ issu de
$x$ contenant $I$ et $J$, et l'unicité de F1 dans $K$ donnerait $I = K_{x,y} =
J$, contre l'hypothèse\footnote{L'argument utilise la forme fermée a) de la
page 22, avec $I \cup J \subset K$ ; F6 tel qu'il est écrit ne met dans $K$
que les intérieurs. Voir la réserve de la section précédente.}.

\emph{Puis $I \cap J = \{x,y\}$.} Si $t \in I \cap J$ n'était pas une
extrémité, il serait intérieur aux deux, et l'on appliquerait ce qui précède
à $I_{x,t}$ et $J_{x,t}$, qui ont les mêmes extrémités et des branches
différentes en $x$ : leurs branches en $t$ différeraient aussi. L'axiome des
lieux réguliers F5 doit alors fournir la contradiction --- le texte s'arrête
sur « on aurait $\mathrm{br}_{t}$ »\footnote{Ce qui manque n'est pas sur la
page, et nous ne le restituons pas. On devine l'intention : en $t$, régulier,
il n'y a que deux branches, et les quatre germes $I_{x,t}$, $I_{t,y}$,
$J_{x,t}$, $J_{t,y}$ doivent s'y répartir d'une manière qu'on espère
contradictoire.}.

Ce qui se dessine, c'est la remarque de la page 2 : deux lieux sont les
extrémités d'au plus deux segments, et de deux exactement quand ils sont sur
un cercle --- deux arcs d'un même cercle, qui ne se rencontrent qu'en leurs
extrémités. C'est le cas circulaire laissé de côté à la page 17. Le dossier
s'arrête là ; le lendemain, 8 juin, le chapitre III (dossier 156-3) reprend
les réseaux « via découpages ».

\end{document}
